% !TEX root = ../main.tex
\chapter{帰納法：リスト・木・再帰処理の仕様}
\label{chap:induction}
\BookIndex{きのうほう}{帰納法}
\BookIndex{List}{\texttt{List}}
\BookIndex{きこうぞう}{木構造}

\begin{chapterabstract}
本章では、帰納法を「任意の大きさのデータ構造に対する仕様確認」として読みます。
学校数学では、帰納法は自然数についての証明として登場することがよくあります。
しかしソフトウェアでも、リスト、木、AST、メニュー構造、コメントスレッド、ディレクトリツリーのような再帰的データを扱うたびに、同じ考え方が現れます。

本章の中心例は、データ移行です。
1件のユーザーを移行する関数があるとき、その関数をリスト全体に適用しても件数が保存されることを証明します。
さらに、木構造全体を変換してもノード数が保存されることを証明します。
ここで使う道具が、Lean の \texttt{induction} です。
\end{chapterabstract}

\begin{learninggoals}
\item 帰納法を「空のケース」と「1つ追加するケース」による全件確認として説明できます。
\item \texttt{List.map}、\texttt{List.length}、\texttt{List.append} に関する基本的な仕様を読めます。
\item Lean の \texttt{induction} が作る \texttt{nil} ケースと \texttt{cons} ケースを読めます。
\item リスト全体のマイグレーションで件数やID列が保存されることを証明できます。
\item 木構造の変換でノード数が保存されることを構造帰納法として理解できます。
\end{learninggoals}

\section{任意の長さや形へ保証を広げる}

実務では、「1件なら正しい」だけでは不十分なことがよくあります。
たとえば、ユーザー1件を新しいスキーマへ移す関数 \texttt{migrateUser} があるとします。
1件の移行がIDを保存することは、比較的簡単に確認できます。
しかし、本番の移行対象は1件とは限りません。
数千件や数百万件の場合もあれば、事前に件数が分からない場合もあります。

ここで確認したい仕様は、任意のユーザー一覧 \texttt{xs} について、
\texttt{xs.map migrateUser} の件数と ID 列が保存されるという主張です。

この「任意の長さ」という条件が重要です。
テストでは、空リスト、1件、2件、大量データなどを用意できます。
しかし、すべての長さを個別にテストすることはできません。
帰納法は、この「すべての長さ」を有限個の証明パターンへ圧縮する道具です。

Lean のリストは、大きく2つの形に分かれます。
空リスト \texttt{[]} と、先頭要素 \texttt{x} に残りのリスト \texttt{xs} を付けた \texttt{x :: xs} です。
この2つの形は、実装上のパターンマッチにもそのまま現れます。

リストに対する帰納法では、\texttt{[]} の場合を基底ケース、\texttt{x :: xs} の場合を帰納ステップと呼びます。
帰納ステップでは、残りのリスト \texttt{xs} について性質が成り立つ、という仮定を使えます。
この仮定を、帰納法の仮定または帰納仮定と呼びます。

\FigurePlaceholder{fig-ch20-list-induction.png}{0.90}{リストの構造帰納法}{fig:ch20-list-induction}

図\ref{fig:ch20-list-induction}では、空リストから始めて、先頭に1つずつ要素を追加する様子を表しています。
帰納法の証明は、この階段を1段ずつ上がれることを示します。
空の段に立てることと、どの段からも次の段へ進めることが分かれば、任意の段に到達できます。

\section{\texttt{List.map}、\texttt{List.length}、\texttt{List.append}}

本章で使うリスト操作は、\texttt{map}、\texttt{length}、\texttt{append} の3つです。
これらは、いずれも日常的なコードに現れます。
Lean では、単なる便利関数ではなく、仕様を書いて証明する対象になります。

\subsection{\texttt{map} は全要素への一括変換}

\texttt{List.map f xs} は、リスト \texttt{xs} の各要素に関数 \texttt{f} を適用します。
ソフトウェアの言葉では、1件の変換をバッチ全体へ持ち上げる操作です。

\subsection{\texttt{length} は件数の仕様を表す}

\texttt{List.length xs} は、リストの長さ、すなわち件数を返します。
実務のマイグレーションでは、件数保存は最初に確認したい性質の1つです。

件数保存とは、変換前後でデータの個数が変わらないという性質です。
リスト \texttt{xs} と変換関数 \texttt{f} について、
\[
  \texttt{(xs.map f).length = xs.length}
\]
が成り立つ、という形で表せます。

この式は、変換結果の中身が正しいことまでは主張していません。
しかし、移行処理が要素を落としたり、重複して増やしたりしないことを確認するための重要な第一歩です。

\subsection{\texttt{append} は分割処理の仕様に出てくる}

\texttt{xs ++ ys} は、リスト \texttt{xs} の後ろに \texttt{ys} をつなぎます。
バッチ処理では、データを分割して処理し、最後に結合することがよくあります。
そのため、\texttt{append} と \texttt{length} の関係は実務上も重要です。

リストを前半 \texttt{xs} と後半 \texttt{ys} に分けた場合、結合後の件数は
\texttt{(xs ++ ys).length = xs.length + ys.length} と表せます。

この性質があると、分割バッチの結果を再結合したときに件数が期待どおりになることを説明しやすくなります。
ただし、実システムではエラー、リトライ、重複排除、トランザクション境界も考える必要があります。
本章では、その中の純粋なリスト変換部分だけを小さく切り出します。

\section{\texttt{induction} の読み方}

Lean の \texttt{induction} は、データ型の形に沿って証明ゴールを分解します。
リストに対して使うと、空リストの場合と、先頭要素を持つ場合に分かれます。

まず、1件の移行関数と、リスト全体への移行関数を定義します。

\begin{listing}[htbp]
\begin{minted}{lean4}
namespace Chapter20

structure UserV1 where
  id : Nat
  name : String
  deriving Repr, BEq

structure UserV2 where
  id : Nat
  displayName : String
  deriving Repr, BEq

def migrateUser (u : UserV1) : UserV2 :=
  { id := u.id, displayName := u.name }

def migrateUsers (xs : List UserV1) : List UserV2 :=
  xs.map migrateUser

end Chapter20
\end{minted}
\caption{\texttt{migrateUser} と \texttt{migrateUsers}}
\label{lst:induction-user-migration}
\end{listing}

ここでは、\texttt{name} というフィールド名を \texttt{displayName} へ変えています。
ただし、IDはそのまま残します。
この例は、後のスキーマ移行の章で扱う大きな例の最小版です。

次に、\texttt{List.map} が長さを保存することを証明します。
この定理は、今後のバッチ移行で繰り返し使う道具になります。

\begin{listing}[htbp]
\begin{minted}{lean4}
theorem map_preserves_length {α β : Type} (f : α → β) :
    ∀ xs : List α, (xs.map f).length = xs.length := by
  intro xs
  induction xs with
  | nil =>
      rfl
  | cons x xs ih =>
      simp [List.map, ih]
\end{minted}
\caption{\texttt{map\_preserves\_length}}
\label{lst:induction-map-length}
\end{listing}

この証明は、日本語では次のように読めます。

\begin{itemize}
\item 任意のリスト \texttt{xs} を取ります。
\item \texttt{xs} について帰納法を使います。
\item \texttt{xs = []} の場合、\texttt{[].map f} も空リストなので、長さはどちらも0です。
\item \texttt{xs = x :: xs} の場合、残りの \texttt{xs} について長さ保存が成り立つと仮定します。
\item 先頭要素にも同じ変換を1回適用するだけなので、リスト全体の長さも保存されます。
\end{itemize}

\texttt{nil} ケースは、空リストです。
\texttt{cons x xs ih} ケースでは、\texttt{x} が先頭要素、\texttt{xs} が残りのリスト、\texttt{ih} が「残りのリストでは定理が成り立つ」という帰納仮定です。

帰納法の仮定 \texttt{ih} は、「任意のリストで成り立つ」という仮定ではありません。
\texttt{cons} ケースの中で、残りのリスト \texttt{xs} について成り立つ、という限定された仮定です。
この範囲を取り違えると、証明の意味を誤解しやすくなります。

\section{サンプル：リスト移行で件数が保存されること}

ここからは、実務に近い形でリスト移行の仕様を証明します。
目標は、\texttt{migrateUsers xs} の件数が \texttt{xs} の件数と等しいことを示すことです。

\subsection{証明方針}

\texttt{migrateUsers} を直接展開すると、\texttt{xs.map migrateUser} になります。
したがって、前節の \texttt{map\_preserves\_length} を使えます。

\begin{listing}[htbp]
\begin{minted}{lean4}
theorem migrateUsers_preserves_length (xs : List UserV1) :
    (migrateUsers xs).length = xs.length := by
  exact map_preserves_length migrateUser xs
\end{minted}
\caption{\texttt{migrateUsers\_preserves\_length}}
\label{lst:induction-migrate-length}
\end{listing}

\texttt{exact} は、\texttt{migrateUsers} を定義展開したゴールに、直前の
\texttt{map\_preserves\_length} を \texttt{migrateUser} へ適用した証明を渡しています。
証明そのものは短くなっています。
これは、\texttt{map\_preserves\_length} という汎用的な道具を先に作っているためです。

\subsection{ID列が保存されること}

件数保存だけでは、移行の正しさとして十分ではありません。
たとえば、件数が同じでも、IDが別の値に書き換わっていたら問題です。
そこで、移行前後のID列が一致することも証明します。

\begin{listing}[htbp]
\begin{minted}{lean4}
def idsV1 (xs : List UserV1) : List Nat :=
  xs.map (fun u => u.id)

def idsV2 (xs : List UserV2) : List Nat :=
  xs.map (fun u => u.id)

theorem migrateUsers_preserves_ids (xs : List UserV1) :
    idsV2 (migrateUsers xs) = idsV1 xs := by
  induction xs with
  | nil =>
      rfl
  | cons x xs ih =>
      simp [migrateUsers, idsV1, idsV2, migrateUser]
\end{minted}
\caption{\texttt{migrateUsers\_preserves\_ids}}
\label{lst:induction-list-ids}
\end{listing}

この証明は、件数保存よりも業務仕様に近いものです。
\texttt{idsV1} は移行前のID一覧、\texttt{idsV2} は移行後のID一覧を取り出します。
定理は、移行してからIDを取り出しても、先にIDを取り出しても同じ列になることを主張しています。

\subsection{\texttt{append} と分割バッチ}

バッチ移行では、一覧を複数のチャンクに分けて処理することがよくあります。
そのとき、前半と後半を結合したリストの長さは、それぞれの長さの和になります。

\begin{listing}[htbp]
\begin{minted}{lean4}
theorem append_length {α : Type} :
    ∀ xs ys : List α, (xs ++ ys).length = xs.length + ys.length := by
  intro xs ys
  induction xs with
  | nil =>
      simp
  | cons x xs ih =>
      simp [ih, Nat.add_assoc, Nat.add_comm]
\end{minted}
\caption{\texttt{append\_length}}
\label{lst:induction-append-length}
\end{listing}

この証明も、\texttt{xs} について帰納法を使っています。
\texttt{ys} について帰納法を使わない点に注意してください。
\texttt{xs ++ ys} の再帰的な定義は、左側の \texttt{xs} の形に沿って進むためです。

帰納法では、どの変数について帰納法を使うかが重要です。
\texttt{xs ++ ys} のように、定義が \texttt{xs} 側を再帰的に見る場合は、まず \texttt{xs} について帰納法を使うと証明しやすくなります。
証明が進まないときは、対象の関数がどの引数を分解しているかを確認してください。

\section{サンプル：木構造変換でノード数が保存されること}

リストは、一方向に伸びるデータ構造です。
実務では、カテゴリツリー、組織図、ファイルツリー、画面部品のツリー、構文木など、より複雑な再帰構造も現れます。
このようなデータでは、空リストと \texttt{cons} ではなく、木の形に沿った帰納法を使います。

木に対する帰納法では、葉のケースと、左右の部分木を持つノードのケースを扱います。
後者では、左右それぞれについて帰納仮定を使えます。

\FigurePlaceholder{fig-ch20-tree-induction.png}{0.90}{木の構造帰納法}{fig:ch20-tree-induction}

図\ref{fig:ch20-tree-induction}では、1つのノードが左部分木と右部分木を持つ様子を表しています。
木全体の性質を示すには、左部分木と右部分木の両方について性質が成り立つことを使います。

\subsection{木構造とノード数}

まず、値を持たない葉 \texttt{leaf} と、左右の部分木と値を持つ \texttt{node} からなる木を定義します。
さらに、木全体に関数を適用する \texttt{Tree.map} と、ノード数を数える \texttt{Tree.size} を定義します。

\begin{listing}[htbp]
\begin{minted}{lean4}
inductive Tree (α : Type) where
  | leaf : Tree α
  | node : Tree α → α → Tree α → Tree α
  deriving Repr

namespace Tree

def map {α β : Type} (f : α → β) : Tree α → Tree β
  | Tree.leaf => Tree.leaf
  | Tree.node l x r => Tree.node (map f l) (f x) (map f r)

def size {α : Type} : Tree α → Nat
  | Tree.leaf => 0
  | Tree.node l _ r => 1 + size l + size r

end Tree
\end{minted}
\caption{\texttt{Tree} と \texttt{Tree.size}}
\label{lst:induction-tree-def}
\end{listing}

ここでは、\texttt{leaf} のサイズを0、\texttt{node} のサイズを「1 + 左部分木のサイズ + 右部分木のサイズ」としています。
この定義は、値の個数を数える仕様です。

\subsection{木の \texttt{map} はノード数を保存する}

木の各ノードに入っている値を変換しても、木の形そのものは変わりません。
したがって、ノード数は保存されるはずです。
この主張を Lean で証明します。

\begin{listing}[htbp]
\begin{minted}{lean4}
theorem tree_map_preserves_size {α β : Type} (f : α → β) :
    ∀ t : Tree α, Tree.size (Tree.map f t) = Tree.size t := by
  intro t
  induction t with
  | leaf =>
      rfl
  | node l x r ihL ihR =>
      simp [Tree.map, Tree.size, ihL, ihR]
\end{minted}
\caption{\texttt{tree\_map\_preserves\_size}}
\label{lst:induction-tree-size}
\end{listing}

この証明では、\texttt{node} ケースで \texttt{ihL} と \texttt{ihR} が現れます。
\texttt{ihL} は左部分木についての帰納仮定であり、\texttt{ihR} は右部分木についての帰納仮定です。
これらを使って左右の部分木のサイズが保存されることを示し、その結果として木全体のサイズ保存を示します。

\subsection{ユーザー木の移行に適用する}

木にユーザー情報が入っている場合も同じです。
\texttt{Tree.map migrateUser} によって各ノードの中身を移行しても、木の形は変わりません。

\begin{listing}[htbp]
\begin{minted}{lean4}
def migrateUserTree (t : Tree UserV1) : Tree UserV2 :=
  Tree.map migrateUser t

theorem migrateUserTree_preserves_size (t : Tree UserV1) :
    Tree.size (migrateUserTree t) = Tree.size t := by
  simpa [migrateUserTree] using tree_map_preserves_size migrateUser t
\end{minted}
\caption{\texttt{migrateUserTree\_preserves\_size}}
\label{lst:induction-migrate-tree}
\end{listing}

リストの場合と同じく、ここでも一般定理を具体的なドメインへ適用しています。
\texttt{tree\_map\_}\allowbreak\texttt{preserves\_size} は、任意の型と任意の変換関数について成り立ちます。
\texttt{migrateUserTree\_}\allowbreak\texttt{preserves\_size} は、それをユーザー移行へ特殊化した定理です。

\section{実務への読み替え}

本章で見た証明は小さいものですが、扱っている考え方は実務でもよく使います。

帰納法で扱いやすい仕様には、次のようなものがあります。
\begin{itemize}
\item リスト移行で件数が保存されます。
\item 全要素についてIDやキーが保存されます。
\item チャンク分割しても、結合後の件数が期待どおりになります。
\item 木構造の変換でノード数や形が保存されます。
\item 再帰的な設定ファイル変換で、キーの構造が保存されます。
\end{itemize}

一方、帰納法だけですべてが解決するわけではありません。
たとえば、移行後の値が外部サービス上でも正しく解釈されるか、DBトランザクションが失敗時に安全にロールバックされるか、分散処理で重複が発生しないか、といった問題には別のモデル化が必要です。

\section{本章のまとめ}

帰納法は、任意の長さのリストや任意の深さの木に対する仕様を扱うための道具です。

リストでは、空リストのケースと、先頭要素を追加する \texttt{cons} ケースを確認します。

\texttt{List.map} の件数保存は、バッチ移行の基本仕様として使えます。

\texttt{append} の長さに関する性質は、分割バッチの説明や結合処理の説明に使えます。

木構造では、左右の部分木に対する帰納仮定を使って、木全体の性質を証明します。

Lean の \texttt{induction} により、有限個の実行例ではなく、任意の大きさのデータ構造について性質を確認できます。
